\documentclass[conference]{IEEEtran}
\usepackage[utf8]{inputenc}
\usepackage[T1]{fontenc}
\usepackage{cite}
\usepackage{amsmath,amssymb}
\usepackage{graphicx}
\usepackage{booktabs}
\usepackage{url}

\title{Evaluating a Deterministic Validator Plus Single‑Iteration Repair Pipeline for LLM‑Generated Network Configurations}

\author{
\IEEEauthorblockN{Autonomous AI Researcher}
\IEEEauthorblockA{CNSM 2026 Agentic AI Researcher Track}
}

\begin{document}

\maketitle

\begin{abstract}
We preregistered and executed a paired confirmatory experiment (N=300 paired intent\ensuremath{\rightarrow}configuration tasks; preregistration id cand-2026-001-A-prereg-v1, preregistration SHA-256 archived) to evaluate whether a minimal guarded pipeline---single‑shot LLM generation (declared model: gpt-5-mini) followed by a frozen deterministic configuration validator and at most one automated repair iteration---reduces first‑pass misconfiguration relative to plain single‑shot generation. Execution completed 600 episodes and used 301 model calls. Observed outcomes: baseline = 299/300 valid (99.667\%); guarded = 300/300 valid (100.0\%); paired counts n11=299, n10=1, n01=0, n00=0. Exact McNemar two‑sided p = 1.00; absolute paired difference = 0.0033333; paired bootstrap 95\% CI = [0.00, 0.01] (10,000 resamples, seed 1729). The single permitted automated repair corrected the sole invalid baseline output (repair success 1/1; Clopper--Pearson 95\% CI = [0.025, 1.0]). The preregistered planning assumption used for power calculations (planned baseline pass rate = 0.40) was contradicted by the observed near‑ceiling baseline, removing statistical headroom for the preregistered >=50\% relative reduction target. We report preregistered design choices, execution accounting, confirmatory statistics, a post‑hoc inferential sensitivity bound tied to the observed contingency counts, and artifact‑grounded recommendations for future NetOps confirmatory evaluations. All run artifacts and analysis outputs are archived in the execution manifest referenced in the Disclosure Statement.
\end{abstract}

\section{Introduction}
Translating operator intent to device configuration is a core NetOps task. Large language models (LLMs) have been proposed as intent\ensuremath{\rightarrow}configuration translators that can accelerate operations, but probabilistic generation risks misordering, omission, and unsafe commands. Recent literature advocates hybrid patterns that pair probabilistic generators with deterministic validators/simulators and bounded repair loops as operational floor‑safety nets. However, rigorous preregistered evaluations that quantify how tightly bounded guarded pipelines change first‑pass misconfiguration rates under controlled sampling semantics are limited.

We implement and preregister a narrowly scoped paired experiment to evaluate the operational effect of a minimal guarded pipeline composed of: (1) a single shared LLM candidate per task (single‑shot generation), (2) a frozen deterministic validator performing canonicalization and rule/dependency checks, and (3) at most one automated repair attempt when the validator reports violations. The core research question is: Can this minimal guarded pipeline reduce first‑pass misconfiguration rate relative to plain single‑shot generation under a controlled paired design? Our contributions are (i) a preregistered paired experiment with deterministic seed generation and archived per‑task artifacts, (ii) confirmatory statistics derived from those artifacts, (iii) a compact post‑hoc sensitivity quantification tied to the observed contingency counts, and (iv) artifact‑grounded recommendations for future NetOps confirmatory evaluations.

\section{Related Work}
Work applying LLMs to network configuration and intent translation has grown rapidly. Recent system demonstrations show promising intent\ensuremath{\rightarrow}configuration synthesis in lab settings but also document failure modes---misordering, omission, and overconfidence---that affect operational safety [10.1002/nem.2313]. Parallel work in adjacent domains (medical QA, code repair) shows that verifier‑assisted and retrieval‑augmented pipelines reduce factual errors and improve grounding, motivating similar guarded architectures for NetOps [10.36227/techrxiv.176784505.56675782/v1, 10.2139/ssrn.6608518]. Classic deterministic network verification (header‑space analysis, specification‑based checking) provides algorithmic foundations for validators that can statically check reachability and policy constraints without live execution [10.1145/3098822.3098834]. Survey and review articles emphasize the generate\ensuremath{\rightarrow}verify\ensuremath{\rightarrow}repair design pattern and highlight auditability and human‑oversight concerns as open challenges for production NetOps adopters [10.6028/nist.ai.100-2e2023]. Our study positions itself at the intersection of these strands by executing a preregistered paired evaluation that directly measures first‑pass validity under an explicit validator+repair floor‑guardrail.

\section{Methodology}
Preregistration, hypotheses, and planning assumptions: The study was preregistered prior to observing outcomes (cand-2026-001-A-prereg-v1). The two preregistered confirmatory hypotheses were: H1 --- the guarded pipeline (deterministic validator + <=1 automated repair) reduces the first‑pass misconfiguration rate by >=50\% (relative) versus plain single‑shot generation; H2 --- one or fewer automated repair iterations recover >=75\% of initially invalid configurations. The preregistration explicitly used a planning assumption of baseline pass rate = 0.40 when computing the sample size and detectable effect: a 50\% relative reduction in the preregistered framing maps the baseline pass rate 0.40 to an expected guarded pass rate \ensuremath{\approx}0.70 (absolute +0.30). Under conservative McNemar approximations reported in the preregistration, N=300 paired tasks provided >80\% power to detect absolute differences on the order of 0.20 at \ensuremath{\alpha}=0.05; these numerical planning choices are reported here because they materially shaped the design and interpretation.

Design overview and pairing semantics: We used a paired within‑task design with a deterministically generated seed list of 300 tasks. For each seed the experiment produced one sampled LLM candidate that was shared between the two conditions (shared\_initial\_candidate semantics): baseline (single‑shot acceptance of the sampled candidate) and guarded (the same sampled candidate passed to the frozen deterministic validator and, if invalid, subjected to the preregistered automated repair adapter with at most one repair attempt). The shared‑candidate pairing was chosen to isolate the incremental effect of deterministic validation and bounded automated repair on a fixed generated plan; this choice mirrors operational workflows where a single proposed plan is reviewed and possibly repaired. The tradeoff is inferential: reusing the same candidate increases within‑pair correlation and reduces the number of informative discordant pairs available to McNemar‑style paired tests compared with an independent‑draws design. The preregistration documented this pairing decision and its expected consequences for statistical information and required discordant counts.

Task bank, inclusion/exclusion, and sampling: The task bank comprised 300 deterministic seeds produced by a frozen task generator with prespecified vendor‑template constraints and intent/workflow patterns. Inclusion required that each generated seed parse to the frozen intent template; seeds that failed the deterministic parsing rule would be deterministically replaced from a preregistered reserve list (the preregistration defines replacement rules). The planned episode count was 600 (two episodes per seed), with the primary estimand computed on the set of complete pairs that satisfied the inclusion rules. The preregistered missingness and retry policy allowed a single retry for provider/API failures; the primary analysis rule ('complete\_pair\_primary') drops pairs only if both calls failed. These procedural choices were set before execution and are described here because they determine how failed calls are treated in the confirmatory test.

Model, generation semantics, and execution budget: The declared LLM in the execution contract was gpt-5-mini. Execution settings declared in the plan include one initial generation call per task (initial\_generation\_calls\_per\_task = 1) and a maximum of one repair attempt per task (maximum\_repair\_calls\_per\_task = 1). The preregistered execution budget and stopping rules were fixed: complete all planned pairs (N=300) unless technical/provider failures exhausted the allowed retries or the model‑call budget. These parameters were chosen to bound the automated repair budget and to align the confirmatory estimand with an operationally plausible, low‑repair‑effort floor guard.

Deterministic validator and automated repair adapter: The frozen deterministic validator (declared identifier: deterministic\_netops\_validator\_v5, validator\_version: 5.0) implements canonicalization to a normalized configuration AST, static rule and dependency checks (for example, required sequences such as shutdown\ensuremath{\rightarrow}modify\ensuremath{\rightarrow}restore), and a binary pass/fail decision under the preregistered scoring rule 'full intent‑constraint validation'. The validator also emits structured violation codes used to classify failure categories (preregistered exploratory categories included reachability, ACL/policy, routing/forwarding, and syntax/ordering mismatches) and to steer the automated repair adapter when an attempt is permitted. The automated repair adapter is bounded: it may issue a single repair prompt transformation of the shared candidate and resubmit the transformed candidate to the validator; if the transformed candidate still fails, no additional automated attempts are made under the preregistered design. The scoring rule required that the validator's canonicalized AST satisfy the task's required command sequence and workflow constraints to be scored as a success.

Primary estimand and inferential procedures: The primary estimand is the paired\_success\_rate\_difference\_guarded\_minus\_baseline (absolute difference in per‑condition success rates computed on complete pairs). The preregistered primary hypothesis test is the exact McNemar two‑sided test at \ensuremath{\alpha}=0.05 applied to the paired pass/fail contingency table (cells n11, n10, n01, n00). We also preregistered paired bootstrap percentile 95\% confidence intervals for the absolute paired difference (10,000 resamples) and prespecified sensitivity analyses: (a) treating failed model calls as failures, and (b) stratified subgroup checks by task difficulty or vendor template. Secondary estimands included repair success conditional on initial failure and per‑category reductions in failure counts; secondary p‑values would be adjusted by Benjamini--Hochberg FDR (q=0.05) when applicable. The preregistration explicitly documented that the planned baseline pass‑rate assumption (0.40) drove the detectable effect calculation and that deviations from this assumption would materially change the study's statistical headroom.

Missingness, retries, and contamination controls: The preregistered missingness plan recorded every failed model call, classified failures into provider/API versus technical failures, and allowed a single retry under the adapter policy. The contamination plan required deterministic exact‑and‑AST equivalence checks between test‑task targets and any frozen transformation/template artifacts; exact duplicates would be excluded and replaced deterministically. The primary analysis used a decontaminated paired set per these rules with prespecified sensitivity analyses including flagged/excluded sets.

Archival and scoring traceability (procedural summary): Per the preregistration and execution contract, every execution artifact needed for confirmatory inference was recorded and archived: per‑task sampled candidate (shared\_initial\_candidate), validator canonicalized configuration, parsed\_command\_count, required\_command\_count, violation codes, repair prompt/response when invoked, and the binary pass/fail scoring. The analysis executor uses these archived structured artifacts to construct the paired contingency table and to run the exact McNemar test and bootstrap CIs. For interpretability of the preregistered power statements we retained the original planning numerics (baseline pass rate = 0.40 \ensuremath{\rightarrow} target guarded \ensuremath{\approx} 0.70) in the public registration so readers understand the mapping between the preregistered relative effect target and the absolute differences the design sought to resolve.

\section{Results}
Execution accounting and artifact provenance (archived): The preregistered run completed all planned episodes: planned episodes = 600, completed episodes = 600, complete analyzed pairs = 300. Model calls consumed = 301 total (300 shared initial generation calls + 1 repair invocation). The archived deterministic reconciliation/provider audit records show no provider or infrastructure failures and no deterministic exclusions for contamination under the preregistered checks. These execution counts and stage summaries are recorded in the archived execution manifest and analysis bundle and are the direct source for all numeric values reported below.

Primary per‑condition outcomes (archived counts): Baseline (single‑shot) success count = 299/300 (99.6667\%); Guarded (validator + <=1 repair) success count = 300/300 (100.0\%). The paired contingency table derived from the archived analysis artifact is: n11 (both success) = 299; n10 (guarded success, baseline fail) = 1; n01 (guarded fail, baseline success) = 0; n00 (both fail) = 0. Missingness summary in the archive reports zero failed or unscorable episodes in either arm under the preregistered missingness treatment.

Confirmatory inferential results (preregistered analyses executed on archived artifacts): The exact McNemar two‑sided test on the observed contingency yields p = 1.00. The observed absolute paired difference (guarded - baseline success rate) = 1/300 \ensuremath{\approx} 0.0033333. Paired bootstrap percentile 95\% CI for the absolute paired difference = [0.00, 0.01] (10,000 resamples, bootstrap seed = 1729), as computed in the archived analysis artifact. Repair outcomes: the single preregistered automated repair invocation succeeded on the sole initially invalid baseline output (repair success 1/1). The exact Clopper--Pearson 95\% CI for a 1/1 repair success is [0.025, 1.0], reflecting the large uncertainty from a single observed repair case.

Why the McNemar p = 1.00 and what the discordant count implies: Exact McNemar inference is driven entirely by the discordant pair counts (n10 and n01). With only one discordant pair observed (n10+n01 = 1), the two‑sided exact test cannot reject the null of no difference; the minimal nonzero discordant total yields an exact two‑sided p of 1.00 under the test implementation used in the preregistered analysis. Equivalently, the discordant total places a hard upper bound on the resolvable absolute paired difference equal to discordants / 300 = 1/300 \ensuremath{\approx} 0.00333. Put practically, the observed near‑ceiling baseline left essentially no statistical headroom to detect the preregistered >=50\% relative reduction (which, given the preregistration's planning baseline = 0.40, mapped to an absolute change on the order of +0.30). Detecting an absolute +0.30 with N=300 would require on the order of tens of discordant pairs (the preregistration noted \textasciitilde{}90 net discordant pairs aligned with the planned effect size under conservative assumptions); the archived discordant count of one is far below that requirement. These limits and numeric bounds are computed from the archived contingency artifact and bootstrap outputs.

Representative validator and diagnostic traces (artifact‑grounded): For each episode the archived validator trace contains structured diagnostics used for scoring: normalized\_configuration, parsed\_command\_count, required\_command\_count, observed\_touched\_field\_count, violation\_codes, validation\_before/after states, and a repair\_applied flag. In the solitary baseline failure the validator trace documents the missing required command sequence and the canonicalized configuration before and after the automatic repair; the post‑repair canonicalized configuration satisfied the full intent constraints per the validator. These per‑task traces (validator decision fields and repair flags) are the primary provenance for the contingency counts and are archived alongside the analysis artifacts.

Sensitivity and interpretive bounds (post‑hoc, artifact‑grounded): Because the observed discordant total is one, the maximal absolute improvement observable in this dataset equals 1/300. The paired bootstrap CI [0.00, 0.01] and the exact McNemar p together reflect that any larger effect size is incompatible with the observed discordant evidence. Operationally, this implies that confirmatory designs aiming to test large relative improvements must either (a) calibrate the task bank to produce a higher baseline failure prevalence or (b) adopt independent‑draws sampling (increasing discordants at the cost of within‑pair control) --- choices that should be incorporated into power calculations. These precise numerical bounds are reproduced in the archived paired\_contingency\_table.csv and paired\_difference\_figure.svg contained in the analysis bundle.

Reproducibility and archived artifacts: All numerical counts, the bootstrap procedure (seed = 1729, resamples = 10,000), the exact‑test implementation, and every per‑task validator trace and repair prompt/response are included in the archived execution manifest and analysis bundle. The archived analysis artifacts (condition summary, paired contingency table, bootstrap outputs, and deterministic reconciliation/provider audit) reproduce the figures and numbers above when loaded with the provided analysis executor. Readers who wish to reproduce the exact inferential outputs can consult the execution manifest and the bundled analysis artifacts referenced in the Disclosure Statement.

Concise summary of the results section: The guarded pipeline corrected the single observed invalid baseline output (artifact‑grounded repair success = 1/1) but, because the baseline pass rate was far higher than the preregistration planning assumption, the run produced only one discordant pair and therefore lacked statistical information to verify the preregistered >=50\% relative reduction. The run nevertheless demonstrates end‑to‑end execution, archival, and artifact‑grounded diagnostics for a generate\ensuremath{\rightarrow}validate\ensuremath{\rightarrow}repair NetOps pipeline and clarifies concrete sample‑size and sampling‑semantic constraints necessary for informative confirmatory inference in future studies.

\section{Discussion}
The guarded pipeline demonstrably remediated the single invalid baseline output observed in this run: the deterministic validator flagged the violation and the single permitted automated repair produced a canonicalized configuration that the validator accepted. This outcome empirically supports the practical utility of a frozen validator plus a tightly bounded repair step as a floor‑safety mechanism that can convert some validator‑detected failures into valid first‑pass configurations.

However, the run also illustrates an important inferential and design caveat: confirmatory power for paired tests depends critically on the number of informative discordant pairs. Our preregistered shared\_initial\_candidate pairing was chosen to model operational workflows that examine and remediate a single generated plan. That design reduces within‑pair heterogeneity and isolates the effect of validation+repair on a given candidate, but it also reduces the expected discordant pair count relative to independent draws per arm. In our execution the baseline pass rate was 299/300, producing only one discordant pair and thereby imposing a strict numerical ceiling on resolvable absolute improvement (1/300 \ensuremath{\approx} 0.00333). Consequently, the preregistered >=50\% relative‑reduction target---formulated under a planning assumption baseline\_pass\_rate = 0.40---became unattainable given the observed prevalences. The archived preregistration explicitly records the baseline\_pass\_rate = 0.40 planning assumption and the analysis artifacts show the realized counts; together they explain why the confirmatory hypothesis could not be supported despite the guarded pipeline correcting the single failure it encountered.

Operationally useful lessons flow directly from these artifact‑grounded facts. First, when baseline error prevalence is uncertain, researchers should design task banks or sampling procedures to produce an expected number of initial failures (or discordant pairs) aligned with the planned detectable effect. For the preregistered aim in this study, the analysis artifacts indicate that on the order of 90 discordant pairs would have been needed to resolve an absolute change comparable with the preregistered target (\ensuremath{\approx}0.30 absolute). In practice, investigators can (a) increase task difficulty, (b) include adversarial or stress‑test augmentations, or (c) adopt stratified sampling by difficulty to concentrate statistical information where it is most informative. These are operational design options rather than claims about guaranteed outcomes; their selection should be preregistered and justified against expected discordant counts.

Second, choice of pairing semantics must be matched to the operational claim. If the goal is to measure whether validator+repair improves the probability that an operator will accept a generated plan when the generator is sampled afresh per condition, independent per‑arm draws increase discordance and statistical information for McNemar‑style inference. If, instead, the workflow examines and remediates a single proposed plan (the shared\_initial\_candidate semantics used here), then evaluation budgets should anticipate reduced discordant counts and compensate via deliberately higher‑difficulty tasks or larger N. The preregistration tradeoff---higher within‑pair control versus greater discordant information---should be made explicit in power computations; our archived power plan assumed baseline\_pass\_rate = 0.40 and thus did not account for the realized near‑ceiling baseline.

Third, repair‑effect estimands are inherently conditional on initial failures and require sufficient counts of those failures to estimate repair effectiveness precisely. In this run the observed repair success was 1/1, which is artifact‑grounded but statistically imprecise (Clopper--Pearson 95\% CI = [0.025, 1.0]). When repair effectiveness is a primary operational claim, studies should be powered to observe an adequate number of initial failures so that conditional repair success has narrow confidence intervals; appropriate sample sizing depends on the target precision and should be preregistered. The deterministic artifacts archived here (per‑task validator traces and repair prompts) exemplify the level of per‑case traceability required to support conditional estimands and subsequent diagnostics.

Finally, construct and external validity remain constrained in this study. Our validator performed deterministic static canonicalization and rule checks rather than full control‑plane or packet‑level emulation; therefore, runtime semantics and emergent control‑plane interactions were not exercised. The task bank is synthetic and vendor‑template coverage was constrained, and only one declared model (gpt-5-mini) and one validator implementation were used. These limitations narrow direct production generalization but do not diminish the methodological point: artifact‑grounded preregistration combined with deterministic validator traces yields reproducible counts that clarify what was actually tested and which inferential questions remain open.

In sum, the empirical artifact shows that small guarded pipelines can correct individual validator‑detected errors, but proving large relative reductions in misconfiguration rate under paired designs requires aligning task difficulty, pairing semantics, and sample size with the planned estimand. We therefore recommend that future confirmatory NetOps evaluations (a) preregister pairing semantics and the expected discordant‑pair count implied by the power calculation, (b) calibrate task difficulty or use adversarial augmentation to achieve a target initial‑failure prevalence when necessary, and (c) archive per‑task validator and repair traces to support transparent conditional estimands and post‑hoc diagnostics. The archived execution and analysis artifacts accompanying this study provide a concrete example of how those recommendations can be operationalized.

\section{Limitations}
\begin{itemize}
\item Synthetic task bank and constrained vendor‑template scope limit external validity to broad production environments.
\item Single LLM (gpt-5-mini) and a single validator implementation restrict generality across model families and validator designs.
\item The preregistered power plan assumed a baseline pass rate = 0.40; the observed baseline (299/300) contradicts that assumption and removed headroom to detect the preregistered >=50\% relative reduction.
\item Sampling semantics (shared\_initial\_candidate) reused the same initial candidate across paired arms, increasing within‑pair correlation and reducing discordant pairs available to McNemar inference.
\item Validation used deterministic configuration/constraint checks rather than full dynamic emulation of runtime semantic effects; actual runtime consequences of configurations were not executed and remain unmeasured.
\item H2 repair‑success evidence is based on a single initially invalid case (1/1 \ensuremath{\rightarrow} 100\%), yielding a wide Clopper--Pearson 95\% CI = [0.025, 1.0]; larger counts of initial failures are required to estimate repair effectiveness precisely.
\end{itemize}

\section{Conclusion}
We executed a preregistered paired experiment (N=300 pairs) comparing single‑shot LLM generation (gpt-5-mini) to a guarded pipeline consisting of a frozen deterministic validator and a single automated repair iteration. Baseline single‑shot generation yielded 299/300 valid outputs and the guarded pipeline yielded 300/300 valid outputs. The single automated repair corrected the lone invalid baseline output, but the observed near‑ceiling baseline contradicted the preregistered planning assumption (baseline pass rate = 0.40) and removed statistical headroom to detect the preregistered >=50\% relative reduction. Future confirmatory NetOps evaluations should (a) calibrate task difficulty to produce sufficient initial failures or target discordant counts, (b) preregister pairing semantics and repair budgets explicitly, and (c) include emulation to measure runtime semantic effects when operational deployment claims are intended. All raw outputs, validator traces, repair traces, the execution manifest, and analysis artifacts are archived and enumerated in the Disclosure Statement to permit reproduction of the reported counts and intervals.

\section*{Disclosure Statement}
This study was executed autonomously after the autonomy lock under the archived master prompt (provenance/master\_prompt.txt; SHA-256: 1872df1e\allowbreak{}1805d2d9\allowbreak{}6940456c\allowbreak{}a016bd66\allowbreak{}5d1d5196\allowbreak{}add77f5a\allowbreak{}cdf1582b\allowbreak{}b39b15ba). Human role: the Research Director selected the topic family (Generative AI and LLMs for NetOps) and approved the master prompt prior to the autonomy lock; no manual human editing of manuscript technical content, figures, tables, or references occurred after the lock. Preregistration: cand-2026-001-A-prereg-v1 (the preregistration record was created before observing final outcomes; preregistration SHA-256 is recorded in the execution manifest). Declared model and tooling: gpt-5-mini (declared\_model\_version), adapter family hosted\_netops\_gvr\_v1, frozen deterministic validator/repair adapter deterministic\_netops\_validator\_v5. Execution accounting (archived): planned episodes = 600, completed episodes = 600, complete pairs = 300, model\_calls\_used = 301 (300 shared initial generation calls + 1 repair invocation), repair-stage invocations = 1. The execution manifest and analysis artifact bundle enumerate all raw model responses, per‑task validator traces (normalized configurations, parsed\_command\_count, required\_command\_count, violation\_codes, repair\_applied flags), repair prompts/responses, execution logs, and analysis artifacts. The archived execution manifest and analysis bundle are the authoritative source for the numeric counts reported in this manuscript. The autonomous pipeline performed literature search, preregistration, execution, analysis, AI‑peer‑review style checks, and manuscript drafting autonomously after the autonomy lock in accordance with the CNSM Agentic AI Researcher track requirements. The preregistered planning assumption used in power calculations (planned baseline pass rate = 0.40) is explicitly reported here because it materially affects interpretation of the confirmatory result.

\begin{thebibliography}{99}
\bibitem{ref1} Nguyen Tu, Sukhyun Nam, James Won‐Ki Hong, "Intent‐Based Network Configuration Using Large Language Models", 2024, 10.1002/nem.2313
\bibitem{ref2} Oghenekeno Hilkiah Okula, ""The Era of AI Agents for Network Engineering": Intent-Aware Auto Remediation for Network Disruption or Failures Using AI Agents, Large Language Models (LLM) and Model Context Protocol (MCP) Mechanisms", 2026, 10.36227/techrxiv.177004277.71795055/v1
\bibitem{ref3} Kunal Dhanda, "MedRAG: Retrieval-Augmented Generation for Medical QA-Comparing Base and RAG-Augmented LLM Performance on Evidence from Peer-Reviewed Clinical Research", 2026, 10.2139/ssrn.6608518
\bibitem{ref4} Goutam Adwant, Manjari Srivastav, "Evidence Coverage Evaluator: A Comprehensive Framework for Assessing Grounding Quality in Retrieval-Augmented Generation Systems", 2026, 10.36227/techrxiv.176784505.56675782/v1
\bibitem{ref5} Ryan Beckett, Aarti Gupta, Ratul Mahajan, David Walker, "A General Approach to Network Configuration Verification", 2017, 10.1145/3098822.3098834
\bibitem{ref6} Apostol Vassilev, Alina Oprea, Alie Jean Fordyce, Hyrum S. Anderson, "Adversarial machine learning :", 2024, 10.6028/nist.ai.100-2e2023
\end{thebibliography}

\end{document}
